\documentclass[11pt]{letter}
\usepackage[margin=2.5cm]{geometry}
\usepackage{hyperref}

\signature{Marco Narcisi}
\address{Independent Researcher\\Castellalto, Abruzzo, Italy\\
\texttt{mnarc123@pm.me}}

\begin{document}

\begin{letter}{The Editor\\Classical and Quantum Gravity}

\opening{Dear Editor,}

We submit for your consideration our manuscript titled
\emph{``Computational search for emergent spacetime geometry from
cellular automata on dynamic graphs''} for publication in
Classical and Quantum Gravity.

This work presents \textsc{Discretum}, a GPU-accelerated computational
framework for the systematic search of local update rules on dynamic
graphs that produce emergent geometric structures.  Using CMA-ES and
genetic algorithm optimisation over a 14-dimensional parameter space,
we identify rules that generate graphs with Hausdorff dimension
approaching~4 from an initially four-dimensional hypercubic lattice.

While our best rules do not achieve the full target of flat
four-dimensional geometry---the spectral dimension is not improved
relative to the bare lattice, and the Ollivier--Ricci curvature
converges to a negative (hyperbolic) value---we provide a detailed
quantitative analysis of why this occurs.  The connectivity collapse
at large system sizes and the absence of spectral crossover are
consistent with the known failure of Euclidean (acausal) dynamical
triangulations, and we identify the imposition of causal structure as
the critical missing ingredient.

We believe this work makes three contributions relevant to the CQG
readership:
\begin{enumerate}
\item A validated, open-source framework for computational exploration
  of discrete quantum gravity models, including GPU-accelerated
  Ollivier--Ricci curvature and ensemble diagnostics.
\item Quantitative constraints on what simple acausal local rules can
  and cannot achieve: correct Hausdorff dimension scaling but anomalous
  spectral dimension, hyperbolic curvature, and connectivity
  instability.
\item A parameter-space analysis identifying topological dynamics
  (edge addition and node merging) as the dominant factors controlling
  emergent geometry, with state transitions playing a negligible role.
\end{enumerate}

The author declares no competing interests and no external
funding was received for this research.

The complete code is publicly available at
\url{https://github.com/mnarc123/discretum} under the MIT
licence.  All results are reproducible from the provided scripts.

\closing{Sincerely,}

\end{letter}
\end{document}
